% =====================================================================
%  4.1 Gjennomsnittlig og momentan vekstfart — teori og oppgaver
% =====================================================================

\del{1}{Gjennomsnittlig vekstfart}

\begin{ordliste}
vekstfart & Hvor fort en størrelse endrer seg, målt per enhet på $x$-aksen. & \mittsprak\\
gjennomsnittlig vekstfart & Hvor mye $f(x)$ endrer seg i snitt per enhet $x$ på et intervall. & \mittsprak\\
sekant & En rett linje gjennom to punkt på en graf. & \mittsprak\\
intervall & Alle tall mellom to grenser, for eksempel $[a, b]$. & \mittsprak\\
\end{ordliste}

\begin{mustkunne}
\begin{minipage}[t]{0.52\linewidth}
\textbf{Gjennomsnittlig vekstfart} fra $x = a$ til $x = b$:
\[ \frac{f(b) - f(a)}{b - a} = \frac{\Delta y}{\Delta x} \]
\textbf{Hvorfor:} Vi deler endringen i $y$ på endringen i $x$. Det er stigningstallet fra 3.1.
Linja gjennom $(a, f(a))$ og $(b, f(b))$ er \emph{sekanten}, og stigningstallet til sekanten er den gjennomsnittlige vekstfarten.\par
\smallskip
\textbf{Enhet:} enheten til $f$ per enhet til $x$, for eksempel km per time.
Negativ vekstfart betyr at $f$ synker i snitt.
\end{minipage}\hfill
\begin{minipage}[t]{0.45\linewidth}
\vspace{0pt}\centering
\begin{tikzpicture}[baseline=(current bounding box.north)]
\begin{axis}[width=\linewidth, height=0.85\linewidth, axis lines=middle, xmin=-0.4, xmax=5.2, ymin=-2.2, ymax=9.5,
  xtick={1,4}, xticklabels={$a$,$b$}, ytick={-1,8}, yticklabels={$f(a)$,$f(b)$},
  tick label style={font=\scriptsize}, xlabel={$x$}, ylabel={$y$}, label style={font=\small}]
\addplot[domain=-0.3:4.45, thick, black!60, samples=40] {x^2 - 2*x};
\addplot[domain=0:4.8, thick, hoved] {3*x - 4};
\draw[dashed, graa] (axis cs:1,-1) -- (axis cs:4,-1) -- (axis cs:4,8);
\addplot[only marks, mark=*, mark size=1.5pt, hoved] coordinates {(1,-1) (4,8)};
\node[font=\scriptsize, anchor=north] at (axis cs:2.5,-1.1) {$\Delta x$};
\node[font=\scriptsize, anchor=west] at (axis cs:4.1,3.5) {$\Delta y$};
\node[font=\scriptsize, hoved, anchor=south east] at (axis cs:2.8,4.6) {sekant};
\end{axis}
\end{tikzpicture}
\end{minipage}
\end{mustkunne}

\begin{eksempel}[1 — fra formel]
La $f(x) = x^2 - 2x$. Finn gjennomsnittlig vekstfart på $[1, 4]$.

\begin{stegtabell}
$f(1) = 1^2 - 2\cdot 1 = -1$ & Funksjonsverdien i starten av intervallet.\\
$f(4) = 4^2 - 2\cdot 4 = 8$ & Funksjonsverdien i slutten av intervallet.\\
$\dfrac{f(4) - f(1)}{4 - 1} = \dfrac{8 - (-1)}{3}$ & Endring i $y$ delt på endring i $x$.\\
$= \dfrac{9}{3} = 3$ & Gjennomsnittlig vekstfart er $3$.\\
Kontroll: & Sekanten $y = 3x - 4$ går gjennom $(1, -1)$ og $(4, 8)$.\\
\end{stegtabell}
\end{eksempel}

\begin{eksempel}[2 — fra tabell, med enhet]
Tabellen viser temperaturen $T$ (i $^\circ$C) utover en dag. Finn gjennomsnittlig vekstfart fra kl. 6 til kl. 12.

\begin{center}
\begin{tabular}{c|ccccc}
kl. & 6 & 9 & 12 & 15 & 18\\ \hline
$T$ ($^\circ$C) & 4 & 8 & 13 & 16 & 12\\
\end{tabular}
\end{center}

\begin{stegtabell}
$\dfrac{T(12) - T(6)}{12 - 6} = \dfrac{13 - 4}{6}$ & Les av $13$ og $4$ i tabellen. Intervallet er 6 timer.\\
$= \dfrac{9}{6} = 1,5$ & Regn ut.\\
$1,5\ {}^\circ\text{C}$ per time & Enhet: $^\circ$C (for $T$) per time (for $x$).\\
Kontroll: & $6 \cdot 1,5 = 9$, og $4 + 9 = 13$. Det stemmer.\\
\end{stegtabell}
\end{eksempel}

\begin{vanligfeil}
\begin{itemize}[leftmargin=1.2em,itemsep=1pt,topsep=0pt]
  \item $\dfrac{f(a) + f(b)}{2}$ er feil. Det er gjennomsnittet av to funksjonsverdier. For $f(x) = x^2 - 2x$ på $[1, 4]$ gir det $3,5$, men riktig vekstfart er $3$.
  \item $\dfrac{f(a) - f(b)}{b - a}$ er feil. Teller og nevner må regnes i samme retning. Her får vi $-3$ i stedet for $3$.
  \item Svaret «$3$» uten enhet er ikke ferdig. I en kontekst skriver du enheten, for eksempel $1,5\ {}^\circ$C per time. Ikke «timer per $^\circ$C».
  \item Gjennomsnittlig vekstfart $0$ betyr ikke at $f$ er konstant. Moteksempel: $f(x) = x^2$ på $[-2, 2]$ gir $0$, men $f(0) = 0$ og $f(2) = 4$.
\end{itemize}
\end{vanligfeil}

\nivaahode{1}

\oppg{1.1}{1}{\uh}
La $g(x) = x^2 + x$. Finn gjennomsnittlig vekstfart på intervallet.
\begin{tasks}(3)
  \task $[0, 3]$
  \task $[1, 3]$
  \task $[-1, 2]$
\end{tasks}

\oppg{1.2}{1}{\uh}
La $h(x) = x^2 - 6x$. Finn gjennomsnittlig vekstfart på $[1, 3]$. Fullfør.
\begin{tasks}(1)
  \task $h(1) =$ \skrivlinje[2cm] \quad og \quad $h(3) =$ \skrivlinje[2cm]
  \task $\dfrac{h(3) - h(1)}{3 - 1} = \dfrac{\rute - (\rute)}{\rute} =$ \skrivlinje[2cm]
\end{tasks}

\oppg{1.3}{1}{\uh}
En bil kjører langs en vei. Tabellen viser hvor langt bilen har kjørt etter $t$ timer.
\begin{center}
\begin{tabular}{c|cccc}
$t$ (timer) & 0 & 1 & 2 & 3\\ \hline
avstand (km) & 0 & 70 & 150 & 210\\
\end{tabular}
\end{center}
\begin{tasks}(1)
  \task Finn gjennomsnittsfarten fra $t = 0$ til $t = 2$. Husk enhet.
  \task Finn gjennomsnittsfarten fra $t = 1$ til $t = 3$.
  \task Finn gjennomsnittsfarten i hver av de tre timene. Hvilken time gikk det raskest?
\end{tasks}

\nivaahode{2}

\oppg{1.4}{2}{\uh}
Grafen viser $f(x) = -x^2 + 4x$. Les av punkt og finn gjennomsnittlig vekstfart på intervallet.
\begin{center}
\begin{tikzpicture}
\begin{axis}[width=6.4cm, height=4.4cm, axis lines=middle, xmin=-0.6, xmax=4.8, ymin=-1, ymax=4.8,
  xtick={0,1,2,3,4}, ytick={1,2,3,4}, grid=both, grid style={black!10},
  tick label style={font=\scriptsize}, xlabel={$x$}, ylabel={$y$}, label style={font=\small}]
\addplot[domain=-0.1:4.1, thick, hoved, samples=40] {-x^2 + 4*x};
\end{axis}
\end{tikzpicture}
\end{center}
\begin{tasks}(3)
  \task $[0, 1]$
  \task $[1, 4]$
  \task $[0, 3]$
\end{tasks}

\oppg{1.5}{2}{\mh}
En bil koster $240\,000$ kr som ny. Verdien etter $t$ år er $V(t) = 240\,000 \cdot 0,8^t$ kr.
Avrund til nærmeste hundre kroner.
\begin{tasks}(1)
  \task Finn gjennomsnittlig vekstfart fra $t = 0$ til $t = 3$. Husk enhet.
  \task Finn gjennomsnittlig vekstfart fra $t = 3$ til $t = 5$.
  \task Svaret i b) er mindre enn svaret i a) uten fortegn. Forklar hvorfor, selv om bilen taper $20\,\%$ hvert år.
\end{tasks}

\oppg{1.6}{2}{\uh}
La $f(x) = x^2 - 3x$. Da er $f(2) = -2$ og $f(6) = 18$. Finn \textbf{den første feilen}. Regn riktig.
\[ \frac{f(6) - f(2)}{6 - 2} \quad\longrightarrow\quad \frac{18 - (-2)}{6 - 2} \quad\longrightarrow\quad \frac{18 - 2}{4} \quad\longrightarrow\quad 4 \]

\nivaahode{3}

\oppg{1.7}{3}{\uh}
Er påstanden alltid, noen ganger eller aldri sann? Begrunn svaret. Anta $a < b$.
\begin{tasks}(1)
  \task For $f(x) = kx + m$ er gjennomsnittlig vekstfart på $[a, b]$ lik $k$.
  \task For $f(x) = x^2$ er gjennomsnittlig vekstfart på $[a, b]$ positiv.
  \task For $f(x) = x^2$ ligger gjennomsnittlig vekstfart på $[a, b]$ mellom $f(a)$ og $f(b)$.
\end{tasks}

\oppg{1.8}{3}{\uh}
Finn det som mangler.
\begin{tasks}(1)
  \task En funksjon har $f(2) = 5$. Gjennomsnittlig vekstfart på $[2, 6]$ er $-3$. Finn $f(6)$.
  \task La $f(x) = x^2 + px$. Gjennomsnittlig vekstfart på $[1, 5]$ er $3$. Finn $p$.
\end{tasks}

\begin{sjekk}
\begin{tasks}(2)
  \task Finn gjennomsnittlig vekstfart for $g(x) = x^2 + 3x$ på $[0, 2]$.
  \task Temperaturen faller fra $12\ {}^\circ$C kl. 14 til $6\ {}^\circ$C kl. 17. Finn gjennomsnittlig vekstfart med enhet.
  \task! Hva forteller en negativ gjennomsnittlig vekstfart om funksjonen?
\end{tasks}
\end{sjekk}

% ---------------------------------------------------------------------

\del{2}{Momentan vekstfart}

\begin{ordliste}
momentan vekstfart & Vekstfarten i ett bestemt punkt, altså akkurat der. & \mittsprak\\
tangent & En rett linje som så vidt berører grafen i ett punkt. & \mittsprak\\
tilnærming & Et tall som er nær det riktige tallet, men ikke helt likt. & \mittsprak\\
nærme seg & Komme stadig nærmere et bestemt tall uten å trenge å bli lik det. & \mittsprak\\
\end{ordliste}

\begin{mustkunne}
\begin{minipage}[t]{0.52\linewidth}
\textbf{Momentan vekstfart} i et punkt er stigningstallet til \emph{tangenten} i punktet.\par
\smallskip
\textbf{Hvorfor sekantene hjelper:} Ett punkt er ikke nok til å regne et stigningstall. Vi bruker to punkt, $x$ og $x + h$:
\[ \frac{f(x + h) - f(x)}{h} \]
Lar vi $h$ bli mindre ($1;\ 0,1;\ 0,01;\ \ldots$), kommer det andre punktet nærmere det første. Sekanten dreier seg mot tangenten, og stigningstallene nærmer seg et tall. Det tallet er den momentane vekstfarten.\par
\smallskip
Vi gir dette tallet et navn i 4.2.
\end{minipage}\hfill
\begin{minipage}[t]{0.45\linewidth}
\vspace{0pt}\centering
\begin{tikzpicture}[baseline=(current bounding box.north)]
\begin{axis}[width=\linewidth, height=0.95\linewidth, axis lines=middle, xmin=-0.3, xmax=4.8, ymin=-1, ymax=9.5,
  xtick={2}, xticklabels={$x$}, ytick=\empty, tick label style={font=\scriptsize}, xlabel={}, ylabel={$y$}, label style={font=\small}]
\addplot[domain=0:4.1, thick, black!60, samples=40] {0.5*x^2};
\addplot[domain=0.3:4.6, thick, hoved] {3*x - 4};
\addplot[domain=0.6:4.4, thick, hoved] {2.5*x - 3};
\addplot[domain=0.3:4.6, thick, dashed, niv3] {2*x - 2};
\addplot[only marks, mark=*, mark size=1.5pt, black] coordinates {(2,2)};
\addplot[only marks, mark=*, mark size=1.2pt, hoved] coordinates {(3,4.5) (4,8)};
\node[font=\scriptsize, niv3, anchor=north west] at (axis cs:2.3,1.4) {tangent};
\node[font=\scriptsize, hoved, anchor=south east] at (axis cs:2.3,7.5) {sekanter};
\end{axis}
\end{tikzpicture}
\end{minipage}
\end{mustkunne}

\begin{eksempel}[1 — mindre og mindre $h$]
La $f(x) = x^2$. Anslå vekstfarten i $x = 3$. Vi regner $\dfrac{f(3 + h) - f(3)}{h}$ for mindre og mindre $h$. Her er $f(3) = 9$.

\begin{center}
\renewcommand{\arraystretch}{1.3}
\begin{tabular}{c|c|c|c}
$h$ & $f(3 + h)$ & $f(3 + h) - 9$ & $\dfrac{f(3 + h) - 9}{h}$\\ \hline
$1$ & $16$ & $7$ & $7$\\
$0,1$ & $9,61$ & $0,61$ & $6,1$\\
$0,01$ & $9,0601$ & $0,0601$ & $6,01$\\
$0,001$ & $9,006001$ & $0,006001$ & $6,001$\\
\end{tabular}
\end{center}

\begin{stegtabell}
$7;\ 6,1;\ 6,01;\ 6,001$ & Tallene blir mindre og nærmer seg $6$.\\
Momentan vekstfart i $x = 3$ er $6$. & Sekantene nærmer seg tangenten når $h \to$ liten.\\
\end{stegtabell}
\end{eksempel}

\begin{eksempel}[2 — lese av en tangent]
Grafen viser $g(x) = \dfrac{x^2}{4}$ med tangent i punktet $(2, 1)$. Finn vekstfarten i $x = 2$ ved å lese av.

\begin{center}
\begin{tikzpicture}
\begin{axis}[width=6cm, height=4cm, axis lines=middle, xmin=-1, xmax=5.2, ymin=-1.6, ymax=4.2,
  xtick={0,1,2,3,4}, ytick={-1,1,2,3}, grid=both, grid style={black!10},
  tick label style={font=\scriptsize}, xlabel={$x$}, ylabel={$y$}, label style={font=\small}]
\addplot[domain=-0.8:4.5, thick, hoved, samples=40] {x^2/4};
\addplot[domain=-0.5:4.6, thick, dashed, niv3] {x - 1};
\addplot[only marks, mark=*, mark size=1.5pt, black] coordinates {(2,1) (0,-1) (4,3)};
\end{axis}
\end{tikzpicture}
\end{center}

\begin{stegtabell}
Punkt på tangenten: $(0, -1)$ og $(4, 3)$ & Velg punkt der tangenten treffer rutenettet.\\
$\dfrac{3 - (-1)}{4 - 0} = \dfrac{4}{4} = 1$ & Stigningstallet til tangenten.\\
Kontroll med $h = 0,1$ og $h = 0,01$: & Tabell gir $1,025$ og $1,0025$. De nærmer seg $1$.\\
\end{stegtabell}
Avlesning fra en tegnet tangent er \emph{unøyaktig}: en liten feil i linja gir en annen verdi. Tabellen gir et sikrere anslag.
\end{eksempel}

\begin{vanligfeil}
\begin{itemize}[leftmargin=1.2em,itemsep=1pt,topsep=0pt]
  \item Momentan vekstfart er ikke funksjonsverdien. For $f(x) = x^2$ er $f(3) = 9$, men vekstfarten i $x = 3$ er $6$. $f(3)$ sier hvor høyt grafen er. Vekstfarten sier hvor bratt den er.
  \item At sekanten med $h = 1$ gir den momentane vekstfarten gjelder ikke generelt. Moteksempel: for $f(x) = x^2$ i $x = 3$ gir $h = 1$ tallet $7$, ikke $6$. Det stemmer bare for en rett linje.
  \item Vekstfarten må ha enhet. Er $s$ i meter og $t$ i sekunder, er enheten m/s, akkurat som for gjennomsnittlig vekstfart.
\end{itemize}
\end{vanligfeil}

\nivaahode{1}

\oppg{2.1}{1}{\uh}
La $f(x) = x^2$. Vi anslår vekstfarten i $x = 1$. Her er $f(1) = 1$. Fullfør tabellen.
\begin{center}
\renewcommand{\arraystretch}{1.4}
\begin{tabular}{c|c|c|c}
$h$ & $f(1 + h)$ & $f(1 + h) - 1$ & $\dfrac{f(1 + h) - 1}{h}$\\ \hline
$1$ & $4$ & $3$ & $3$\\
$0,1$ & \rute & \rute & \rute\\
$0,01$ & \rute & \rute & \rute\\
\end{tabular}
\end{center}
\begin{tasks}(1)
  \task Hvilket tall nærmer stigningstallene seg?
\end{tasks}

\oppg{2.2}{1}{\uh}
Grafen viser $f(x) = x^2 - 4x + 5$ og tangenten i punktet $(3, 2)$.
\begin{center}
\begin{tikzpicture}
\begin{axis}[width=6cm, height=4.2cm, axis lines=middle, xmin=-0.5, xmax=5.3, ymin=-0.8, ymax=5.6,
  xtick={0,1,2,3,4,5}, ytick={1,2,3,4,5}, grid=both, grid style={black!10},
  tick label style={font=\scriptsize}, xlabel={$x$}, ylabel={$y$}, label style={font=\small}]
\addplot[domain=0.2:4.4, thick, hoved, samples=40] {x^2 - 4*x + 5};
\addplot[domain=1.7:4.4, thick, dashed, niv3] {2*x - 4};
\addplot[only marks, mark=*, mark size=1.5pt, black] coordinates {(3,2) (2,0) (4,4) (2,1)};
\end{axis}
\end{tikzpicture}
\end{center}
\begin{tasks}(1)
  \task Tangenten går gjennom $(2, 0)$ og $(4, 4)$. Finn momentan vekstfart i $x = 3$.
  \task Er funksjonen voksende eller avtagende i $x = 3$? Hvordan ser du det?
  \task Hva er momentan vekstfart i bunnpunktet $(2, 1)$? Hvordan ser tangenten ut der?
\end{tasks}

\oppg{2.3}{1}{\uh}
Avgjør om vekstfarten er gjennomsnittlig eller momentan. Finn størrelsen med enhet når du kan.
\begin{tasks}(1)
  \task Fartsmåleren i bilen viser $70$ km/h akkurat nå.
  \task Bilen kjørte $270$ km på $3$ timer.
  \task Temperaturen var $4\ {}^\circ$C kl. 6 og $10\ {}^\circ$C kl. 9.
  \task Termometeret stiger akkurat nå med $3\ {}^\circ$C per time.
\end{tasks}

\nivaahode{2}

\oppg{2.4}{2}{\mh}
La $f(x) = x^3$.
\begin{tasks}(1)
  \task Regn $\dfrac{f(1 + h) - f(1)}{h}$ for $h = 1;\ 0,1;\ 0,01;\ 0,001$. Lag en tabell.
  \task Anslå vekstfarten i $x = 1$.
  \task Regn $\dfrac{f(2 + h) - f(2)}{h}$ for $h = 0,01$ og $h = 0,001$. Anslå vekstfarten i $x = 2$.
\end{tasks}

\oppg{2.5}{2}{\mh}
En syklist kjører en etappe. Tabellen viser hvor langt hun har kommet etter $t$ sekunder.
\begin{center}
\begin{tabular}{c|ccccccc}
$t$ (s) & 0 & 10 & 20 & 30 & 40 & 50 & 60\\ \hline
$s$ (m) & 0 & 60 & 150 & 270 & 400 & 520 & 620\\
\end{tabular}
\end{center}
\begin{tasks}(1)
  \task Finn gjennomsnittsfarten fra $t = 0$ til $t = 60$. Avrund til én desimal.
  \task Anslå farten ved $t = 30$ ved å regne på $[20, 30]$ og $[30, 40]$.
  \task Hva er forskjellen i betydning mellom svaret i a) og svaret i b)?
\end{tasks}

\oppg{2.6}{2}{\uh}
Vi anslår vekstfarten til $f(x) = x^2 - x$ i $x = 2$ med $h = 0,1$. Finn \textbf{den første feilen}. Regn riktig.
\[ \frac{f(2,1) - f(2)}{2,1 - 2} \ \longrightarrow\ \frac{(4,41 - 2,1) - 2}{0,1} \ \longrightarrow\ \frac{2,31 - 2}{0,1} \ \longrightarrow\ \frac{0,31}{0,1} = 0,031 \]

\nivaahode{3}

\oppg{2.7}{3}{\mh}
La $f(x) = x^2$. Vi undersøker et mønster.
\begin{tasks}(1)
  \task Finn gjennomsnittlig vekstfart på $[2, 8]$.
  \task Anslå vekstfarten i $x = 5$, midt i intervallet. Bruk $h = 0,1$ og $h = 0,01$.
  \task Sammenlign a) og b). Velg selv et nytt intervall, og undersøk om det samme skjer. Hva tror du, og hvorfor kan du ikke vite det sikkert ennå?
\end{tasks}

\oppg{2.8}{3}{\mh}
For en funksjon $f$ er sekantens stigningstall mellom $x = 2$ og $x = 2 + h$ lik $5$ for $h = 1$, $4,1$ for $h = 0,1$ og $4,01$ for $h = 0,01$.
\begin{tasks}(1)
  \task Anslå vekstfarten i $x = 2$.
  \task Kan vekstfarten i $x = 2$ være $4,5$? Begrunn.
  \task Finn en enkel funksjon som passer med tabellen. Kontroller med regning.
  \hint{Regn sekantens stigningstall for noen enkle funksjoner i $x = 2$.}
\end{tasks}

\begin{sjekk}
\begin{tasks}(2)
  \task La $f(x) = x^2 + 1$. Regn $\dfrac{f(0 + h) - f(0)}{h}$ for $h = 0,1$.
  \task En tangent går gjennom $(0, 1)$ og $(2, 5)$. Hva er vekstfarten i berøringspunktet?
  \task! Hvorfor nærmer sekanten seg tangenten når $h$ blir mindre? Forklar kort.
\end{tasks}
\end{sjekk}
